\section{Results and Discussion}
\label{sec:results}

\subsection{LCA Comparison}

\todo{Present the three-way LCA comparison (Reference vs.\ DfD Concept vs.\ Alternative Concept) across all impact categories. Use bar charts or tables.}

\begin{table}[H]
\centering
\caption{LCA results comparison across three design concepts.}
\label{tab:lca_results}
\begin{tabular}{@{}lccc@{}}
\toprule
\textbf{Impact Category} & \textbf{Reference} & \textbf{DfD Concept} & \textbf{Alternative} \\
\midrule
GWP (kg CO$_2$-eq)                & \todo{X} & \todo{X} & \todo{X} \\
Acidification (kg SO$_2$-eq)      & \todo{X} & \todo{X} & \todo{X} \\
Eutrophication (kg PO$_4$-eq)     & \todo{X} & \todo{X} & \todo{X} \\
Resource Depletion (kg Sb-eq)      & \todo{X} & \todo{X} & \todo{X} \\
Human Toxicity (kg 1,4-DCB-eq)    & \todo{X} & \todo{X} & \todo{X} \\
\bottomrule
\end{tabular}
\end{table}

% \begin{figure}[H]
%     \centering
%     \includegraphics[width=0.9\textwidth]{figures/lca_results_bar.pdf}
%     \caption{Comparative LCA results across all three design concepts.}
%     \label{fig:lca_results}
% \end{figure}

\subsection{Disassembly Performance}

\begin{table}[H]
\centering
\caption{Disassembly performance metrics.}
\label{tab:disassembly_results}
\begin{tabular}{@{}lccc@{}}
\toprule
\textbf{Metric} & \textbf{Reference} & \textbf{DfD Concept} & \textbf{Alternative} \\
\midrule
Battery removal steps     & $\sim$10           & 1          & 1 \\
Battery removal time      & 4--5 min           & $<$5 sec   & $<$5 sec \\
Full disassembly steps    & $\sim$10           & $\leq$5    & $\leq$5 \\
Full disassembly time     & 4--5 min           & $<$30 sec  & $<$30 sec \\
Tools required            & X-acto knife       & None       & None \\
Destructive?              & Yes                & No         & No \\
\bottomrule
\end{tabular}
\end{table}

\subsection{Material Recovery Analysis}

\begin{table}[H]
\centering
\caption{Estimated material recovery rates by design concept.}
\label{tab:recovery}
\begin{tabular}{@{}lccc@{}}
\toprule
\textbf{Material} & \textbf{Reference} & \textbf{DfD Concept} & \textbf{Alternative} \\
\midrule
Lithium          & $\sim$0\% & $>$90\% & $>$95\% (reuse + recycle) \\
Cobalt/Nickel    & $\sim$0\% & $>$90\% & $>$95\% \\
Copper           & $\sim$0\% & $\sim$80\% & $\sim$85\% \\
Aluminum         & $\sim$0\% & $\sim$85\% & N/A (PP housing) \\
Plastics         & $\sim$0\% & $\sim$50\% & $>$80\% (mono-material) \\
\bottomrule
\end{tabular}
\end{table}

\subsection{Cost Implications}

\todo{Analyze the cost implications of the redesign:}
\begin{itemize}
    \item \textbf{Additional manufacturing cost}: Snap-fit tooling, battery door mechanism, spring contacts, material marking. Estimate: \todo{\$X per unit increase.}
    \item \textbf{Material savings}: Battery reuse reduces per-device battery cost. If battery is reused 10 times, battery cost per device drops by $\sim$90\%.
    \item \textbf{Regulatory compliance}: DfD design preemptively meets EU Batteries Regulation 2027 requirements, avoiding future redesign costs.
    \item \textbf{Waste management savings}: Reduced fire risk and simplified recycling lower municipal waste processing costs.
\end{itemize}

\subsection{Discussion}

\todo{Interpret the results. Consider:}
\begin{itemize}
    \item Which impact categories show the greatest reduction from the alternative concept?
    \item How does the battery reuse assumption (number of cycles) affect the results?
    \item What are the barriers to adoption? (manufacturer resistance, consumer behavior, cost, standardization governance)
    \item How does this compare to the outright ban approach (UK, Australia)?
    \item What role does the EU Batteries Regulation play as a forcing function?
    \item Limitations: assumptions in LCA modeling, idealized recycling rates, consumer participation uncertainty
\end{itemize}
